\section{引言}\label{sec:intro}

\subsection{研究背景与统计问题}\label{sec:intro-bg}

时间序列的结构变化既可能表现为均值、方差等低阶统计量的改变，也可能表现为局部轨迹几何、周期结构或状态空间形状的改变。滑动窗口和时延嵌入为后一类信号提供了局部化表示：每个时间片段被转换为点云，持续同调进一步从多尺度上提取连通分支和环等结构信息。相关理论奠定了滑动窗口持久性分析的基础\cite{PereaHarer2015}，持续图像等有限维表示则使持续图能够进入统计扫描和机器学习程序\cite{Adams2017}。近年来，拓扑摘要已被用于离线变点检测、在线监测及依赖序列中的结构异常识别\cite{IslambekovYuvarajGel2020,ZhengMakXieXie2023,ChazalLevrardRoyer2024,NieYue2026}。

为提高候选变点的扫描密度，实际分析常取步长 $s$ 小于窗口宽度 $W$。此时，相邻窗口共享原始观测，窗口特征序列由此形成局部依赖。设第 $i$ 个窗口经固定映射得到特征 $Y_i$。对全部窗口顺序实施置换时，有限样本精确性通常依赖联合分布对排列群保持不变，即
\begin{equation}\label{eq:intro-exchangeability}
 (Y_1,\ldots,Y_M)\overset{d}{=}(Y_{\pi(1)},\ldots,Y_{\pi(M)}),
 \qquad \forall\,\pi\in\mathfrak S_M.
\end{equation}
独立同分布是满足该条件的一种充分情形，排列操作直接要求的条件则是窗口特征层面的联合交换性。平稳性、边际同分布以及原始序列的独立同分布性质，在窗口重叠时均不足以一般性地推出式\eqref{eq:intro-exchangeability}。由此形成本文关注的统计错配：科学零假设定义在原始过程的“无变化”层面，置换校准所需的对称性定义在派生窗口特征的联合分布层面。依赖时间序列中的置换理论同样表明，未经修正的排列参考分布可能无法维持名义检验水平\cite{RomanoTirlea2022}。

\subsection{相关工作与研究缺口}\label{sec:intro-lit}

与本文相关的研究可概括为三条脉络。第一类工作研究滑动窗口嵌入、持续同调及持续图向量化，为局部时间片段的拓扑表示提供方法基础\cite{PereaHarer2015,Adams2017}。第二类工作将拓扑摘要引入离线或在线变化检测，主要考察表示对结构变化的辨识能力及特定检测程序的性能\cite{IslambekovYuvarajGel2020,ZhengMakXieXie2023,ChazalLevrardRoyer2024,NieYue2026}。第三类工作研究依赖数据下的置换、dependent wild bootstrap 与依赖乘子自助法，说明时间依赖会改变随机化参考分布，并为若干均值、光滑函数或经验过程问题建立渐近理论\cite{RomanoTirlea2022,Shao2010,BucherKojadinovic2016}。

这些结果尚未直接覆盖本文考察的复合管线：重叠局部窗口依次经过非线性点云映射、持续同调向量化和高维最大扫描后，朴素置换会产生多大程度的有限样本尺寸失真；该失真如何随重叠程度、原序列依赖、观测预算和表示类型变化；依赖感知校准能够在哪些条件下减弱失真，又会在哪些条件下表现为保守或残留反保守。本文将条件组合 $\theta$ 下的有效检验水平记为 $\alpha_{\mathrm{eff}}(\theta)$，并将映射 $\theta\mapsto\alpha_{\mathrm{eff}}(\theta)$ 称为有限样本校准表面。预设设计域内，从实用校准区转入保守区或反保守区的逐格转换构成本文所称的经验校准边界；该边界属于有限样本、点态和设计域内的刻画，不提供全参数空间的统一保证。

\subsection{研究目标与贡献}\label{sec:intro-contrib}

围绕上述缺口，本文形成三项相互衔接的贡献。
\begin{enumerate}
\item \textbf{重叠窗口下交换性失效的机制刻画。} 对原始独立同分布序列和固定局部映射，本文给出窗口特征有限依赖范围的上界；进一步以高斯窗口均值构造精确的非交换性反例，并推导有限样本方差比和长期协方差错配，从而明确原始层面的独立同分布性质为何不能直接传递为窗口层面的排列不变性。
\item \textbf{双预算制度下的有限样本校准表面。} 本文在固定窗口数与固定观测长度两种预算制度、28 个预先设定的零假设配置上，比较朴素置换与依赖乘子校准，系统刻画重叠程度、原序列依赖、预算制度和表示共同作用下的 I 类错误率，并以预先冻结的主对比量化跨预算制度的失校准幅度。
\item \textbf{表示敏感性与校准有效性的分层评价。} 本文利用匹配零假设经验临界值构造尺寸标准化功效，比较拓扑表示与矩表示在形状变化和均值平移下的信号选择性。该设计将表示所保留的信息与参考分布的松紧程度区分开来，同时保留二者在有限样本下可能发生交互的空间。
\end{enumerate}

全文其余部分安排如下：第~\ref{sec:method} 节定义重叠窗口特征、两类表示、最大扫描统计量及校准方法；第~\ref{sec:theory} 节分析有限依赖、交换性失效与方差机制；第~\ref{sec:protocol} 节说明评价对象和实验设计；第~\ref{sec:h0} 节报告零假设下的有限样本校准结果；第~\ref{sec:h1} 节分析备择下的检出与定位性能；第~\ref{sec:discussion} 节讨论方法含义、适用范围与局限；第~\ref{sec:conclusion} 节总结全文。
